\section{The normalized quantum-group comparison and the compact limit}
\label{sec:normalizedsixj}

The further comparison has two distinct outcomes. The finite
screening operator can be matched explicitly to a specified
$b$-$6j$ normalization: its double-Gamma factors reduce to
ordinary Gamma factors and its analytic quantum-group matrix
is a sine matrix. Here the operator is fixed by the independent
period calculation in Section~\ref{sec:screenedfusion}; its
interpretation through a continued infinite spectral representation
retains the analytic input specified there. Its auxiliary angle Gram matrix
approaches rank two, whereas a nondegenerate shrinking compact
tetrahedron approaches rank three. A direct calculation for
the regular compact family gives a cubic classical saddle
and fixes the proposed joint scale. The exact integration coordinate
requires care: a previously stated pole accumulation at the saddle
came from an incorrect shift. The corrected dictionary below and
the local pole-free strip identify the candidate quantum fold.
Section~\ref{sec:quantumAiry} proves the uniform quantum Airy formula
for this pole-cleared regular family and bounded real transition
parameters, using the exact-integral contour and tail analysis.
The finite screening calculation and the classical fold here do not
supply those additional estimates.

\subsection{Fixing the Teschner--Vartanov normalization}

Normalization changes coefficients because it changes the basis
vectors being compared. If $f=M g$, and we replace $f_j$ by
$f_j/a_j$ and $g_k$ by $g_k/b_k$, the new matrix has entries
$(b_k/a_j)M_{jk}$. This elementary rule accounts for the ratios
of trivalent factors below. A trivalent factor belongs to one
three-label coupling; a channel contains two such couplings.
The Plancherel density is the weight in the continuous analogue
of an orthonormal basis expansion.

Write $R^{\mathrm{an}}$ for the analytic $b$-$6j$ symbol and
$R$ for the symmetric symbol in~\cite{TeschnerVartanov}.
Their equations (4.5)--(4.6) give the standard Virasoro fusion
kernel as $F_{st}=(N_s/N_t)m(t)R^{\mathrm{an}}_{st}$, where
\begin{equation}
 N(z,y,x)=
 \frac{\Gamma_b(2Q-2z)\Gamma_b(2y)\Gamma_b(2x)}
 {\Gamma_b(2Q-x-y-z)\Gamma_b(Q-x-y+z)
  \Gamma_b(x+z-y)\Gamma_b(y+z-x)}.
 \label{eq:TVnormalization}
\end{equation}
Here $N_s$ and $N_t$ are products of the two trivalent
normalizations in each channel. On the principal contour,
$m(a)=|S_b(2a)|^2$ is the Plancherel density. The trivalent
factor distinguishing the symmetric symbol is
\[
 D(z,y,x)=
 \bigl[S_b(2Q-x-y-z)S_b(Q-x-y+z)
 S_b(x+z-y)S_b(y+z-x)\bigr]^{-1/2}.
\]
If $D_s,D_t$ are its channel products, then
$R=(D_s/D_t)R^{\mathrm{an}}$. Equivalently, equations
(4.10)--(4.11) of that paper read
\begin{equation}
 G_{st}=\sqrt{\frac{C_s}{C_t}}\,F_{st}
       =\sqrt{m(s)m(t)}\,R_{st},
 \label{eq:TVunitarykernel}
\end{equation}
where $C_s,C_t$ are the corresponding products of DOZZ
three-point functions. These identities specify the scalar
factors that were left open in Section~\ref{sec:screenedfusion}.

For our ordering of fields, use TV's external labels
\begin{equation}
 (p,q,r,d_{\mathrm{out}})
   =(\alpha_1,\alpha_3,\alpha_2,Q-\alpha_4).
 \label{eq:TVreflectedlabels}
\end{equation}
Reflecting the infinity charge preserves its Virasoro weight.
At the one-screening resonance,
\[
 d_{\mathrm{out}}=p+q+r+b,\qquad
 s_j=p+q+jb,\qquad t_k=q+r+kb,\qquad j,k\in\{0,1\}.
\]
Consequently each channel contains one zero-screening and one
one-screening trivalent vertex. Explicitly,
\[
 N_{s,j}=N(s_j,q,p)N(d_{\mathrm{out}},r,s_j),\qquad
 N_{t,k}=N(t_k,r,q)N(d_{\mathrm{out}},t_k,p).
\]
The reflection convention matters for the analytic symbol,
and is retained throughout the following formulas. Initially
the weights and auxiliary factors are generic, with
$\delta_0\delta_1\ne0$; balanced limits are taken afterward.

\subsection{The double-Gamma factors and the finite sine matrix}

The shift identity
\[
 \frac{\Gamma_b(x+b)}{\Gamma_b(x)}
       =\frac{\sqrt{2\pi}\,b^{bx-1/2}}{\Gamma(bx)}
\]
immediately gives
\begin{align}
 N(x+y,y,x)&=\Gamma_b(Q)^{-1},\notag\\
 N(x+y+b,y,x)&=
 \frac{\Gamma(1+b^2)\Gamma(2bx)\Gamma(2by)
                 \Gamma(2-2b(x+y))}
      {(2\pi)^2b^{1+b^2}\Gamma_b(Q)}.
 \label{eq:TVscreeningN}
\end{align}
Thus
\[
 N_{s,j}=\mathcal K_b n_{s,j},\qquad
 N_{t,k}=\mathcal K_b n_{t,k},\qquad
 \mathcal K_b=
 \frac{\Gamma(1+b^2)}
      {(2\pi)^2b^{1+b^2}\Gamma_b(Q)^2},
\]
with
\begin{align}
 n_{s,0}&=\Gamma(1-\delta_0)\Gamma(\mu_2)\Gamma(\mu_4),
 &n_{s,1}&=\Gamma(1+\delta_0)\Gamma(\mu_1)\Gamma(\mu_3),
 \notag\\
 n_{t,0}&=\Gamma(1-\delta_1)\Gamma(\mu_1)\Gamma(\mu_4),
 &n_{t,1}&=\Gamma(1+\delta_1)\Gamma(\mu_2)\Gamma(\mu_3).
 \label{eq:TVfiniteN}
\end{align}
All $b$-dependence in these normalization products is in the
same scalar $\mathcal K_b$.

Rescale the screened channel vectors by
$e_{s,j}=\mathcal F_{s,j}/N_{s,j}$ and
$e_{t,k}=\mathcal F_{t,k}/N_{t,k}$. Their matrix is
$\mathsf B_{jk}=(n_{t,k}/n_{s,j})\mathsf M_{jk}$, where
$\mathsf M$ is~\eqref{eq:screeningfusionmatrix}. Gamma reflection
gives the explicit answer
\begin{equation}
 \boxed{\quad
 \mathsf B=
 \frac1{\sin(\pi\delta_1)}
 \begin{pmatrix}
 \sin(\pi\mu_2)&-\sin(\pi\mu_4)\\
 \sin(\pi\mu_3)&-\sin(\pi\mu_1)
 \end{pmatrix}.
 \quad}
 \label{eq:analytic6jscreening}
\end{equation}
For example its upper-left entry is
\[
 \frac{n_{t,0}}{n_{s,0}}\mathsf M_{00}
 =\frac{\Gamma(1-\delta_1)\Gamma(\delta_1)}
        {\Gamma(\mu_2)\Gamma(1-\mu_2)}
 =\frac{\sin(\pi\mu_2)}{\sin(\pi\delta_1)}.
\]
This is the induced finite action of $m(t)R^{\mathrm{an}}$
in the specified convention, understood through the exact
normalization identity on the screened period space. Identifying
it with the residue of the full infinite spectral operator
retains the continuation input of Section~\ref{sec:screenedfusion};
no pointwise substitution into a singular Barnes integral is
being made. Its determinant is
$-\sin(\pi\delta_0)/\sin(\pi\delta_1)$.

For a nonbalanced example, take
$\mu=(1/4,1/3,1/2,11/12)$. These weights lie in $(0,1)$ and
sum to two, with $\delta_0=1/4$ and $\delta_1=1/6$.
The finite analytic matrix is then
\[
 \mathsf B=
 \begin{pmatrix}
 \sqrt3&-(\sqrt6-\sqrt2)/2\\
 2&-\sqrt2
 \end{pmatrix},\qquad \det\mathsf B=-\sqrt2.
\]
The determinant agrees with
$-\sin(\pi/4)/\sin(\pi/6)$. Its entries need not have modulus
at most one: the preserved pairing is indefinite, as the next
formula makes explicit, rather than a probability inner product.

For reference, put
$\sigma_i=\sin(\pi\mu_i)$ and
$d_\nu=\sin(\pi\delta_\nu)$. After removing the common
$E_b/\mathcal K_b$ from these vectors, their Hermitian matrices
are
\begin{equation}
 H_s^{\mathrm{an}}=
 \frac{\pi^4}{d_0}
 \diag\left(\frac1{\sigma_2\sigma_4},
            -\frac1{\sigma_1\sigma_3}\right),\qquad
 H_t^{\mathrm{an}}=
 \frac{\pi^4}{d_1}
 \diag\left(\frac1{\sigma_1\sigma_4},
            -\frac1{\sigma_2\sigma_3}\right).
 \label{eq:analytic6jpairing}
\end{equation}
They follow from $H_s^{\mathrm{an}}=\diag(h_jn_{s,j}^2)$,
and satisfy $\mathsf B^*H_s^{\mathrm{an}}\mathsf B
=H_t^{\mathrm{an}}$.

\subsection{Symmetric normalization and the residue prescription}

The square roots in~\eqref{eq:TVunitarykernel} also have a
concrete interpretation at a pinch. Write
$\alpha_s=s_j+\varepsilon\xi$, with fixed
$\xi\ne0,-1$. The product of three-point factors has the form
\[
 C_s\sim\frac{\Lambda_*h_{s,j}}
                {\varepsilon^2\xi(\xi+1)}.
\]
The common $\Lambda_*$ includes the infinity reflection factor
and cancels in channel ratios. With square-root branches fixed
by continuation, the DOZZ-normalized vector
$\mathcal G_s=\sqrt{C_s}\mathcal F_s$ obeys
\begin{equation}
 \lim_{\varepsilon\to0}
 \varepsilon\sqrt{\frac{\xi(\xi+1)}{\Lambda_*}}\,
 \mathcal G_s(s_j+\varepsilon\xi;\varepsilon)
       =\sqrt{h_{s,j}}\mathcal F_{s,j}.
 \label{eq:regularizedTVvector}
\end{equation}
Hence its finite analytic matrix is
\[
 \mathsf O=\diag(\sqrt{h_s})\mathsf M
                 \diag(\sqrt{h_t})^{-1},
 \qquad \mathsf O^T\mathsf O=I.
\]
This is complex orthogonality: negative Hermitian coefficients
give factors of $i$ in analytic square roots.

Using absolute square roots instead gives a real matrix with
an indefinite pairing:
\begin{equation}
 \boxed{\quad
 \mathsf U=
 \frac{\operatorname{sgn}(d_1)}{\sqrt{|d_0d_1|}}
 \begin{pmatrix}
 \sqrt{\sigma_1\sigma_2}&-\sqrt{\sigma_3\sigma_4}\\
 \sqrt{\sigma_3\sigma_4}&-\sqrt{\sigma_1\sigma_2}
 \end{pmatrix}.
 \quad}
 \label{eq:symmetric6jscreening}
\end{equation}
Indeed,
\[
 \sigma_1\sigma_2-\sigma_3\sigma_4=d_0d_1,\qquad
 \mathsf U^T J_s\mathsf U=J_t.
\]
\[
 J_s=\operatorname{sgn}(d_0)\diag(1,-1),\quad
 J_t=\operatorname{sgn}(d_1)\diag(1,-1).
\]
Equations~\eqref{eq:analytic6jscreening} and
\eqref{eq:symmetric6jscreening} are fully normalized finite
operator comparisons. The logarithmic basis of
Section~\ref{sec:screenedfusion} supplies the balanced limits;
one cannot set $d_0=d_1=0$ in the displayed matrices.

There is a separate reason to retain the operator prescription.
The analytic continuation of the Plancherel factor is
\begin{align}
 m(a)&=S_b(2a)S_b(2Q-2a)\notag\\
 &=-4\sin\bigl(\pi b(2a-Q)\bigr)
        \sin\bigl(\pi(2a-Q)/b\bigr).
 \label{eq:continuedPlancherel}
\end{align}
The two $S_b$ shift equations prove the second equality.
At the screened labels this becomes
\begin{equation}
 m(s_j)=
 4\sin\bigl(\pi[\delta_0-(2j-1)b^2]\bigr)
       \sin(\pi\delta_0/b^2),
 \label{eq:screenedPlancherel}
\end{equation}
with the analogous formula for $t_k,\delta_1$. These factors
oscillate for small real $b$ and vanish at balanced weights.
The trivalent symmetric normalizations also encounter
$S_b(Q+jb)=0$. Thus dividing a finite matrix by pointwise
Plancherel values does not define a bare-symbol limit.

\subsection{The actual screened Gram matrix has rank two at the limit}

The tetrahedral dictionary of~\cite[Section~3]{TeschnerVartanov}
is
\[
 A_e=\exp(\pi i-2\pi ib\alpha_e)=e^{i\eta_e}.
\]
The vertex triples are $(1,2,s),(3,4,s),(1,4,t),(2,3,t)$;
compare the original convention in~\cite{MurakamiYano}.
Here $A_e$ is an exponentiated angle, distinct from surface area.
Put
\[
 \vartheta_i=\pi\mu_i,\qquad h=2\pi b^2,\qquad
 \sum_i\vartheta_i=2\pi.
\]
With the reflected labels~\eqref{eq:TVreflectedlabels}, a
face ordering gives the formal angle Gram matrix
\begin{equation}
 \mathbf G_{jk}(h)=
 \begin{pmatrix}
 1&\cos(\vartheta_2+\vartheta_3+kh)&\cos\vartheta_2&
                                      \cos(\vartheta_4-h)\\
 \cos(\vartheta_2+\vartheta_3+kh)&1&\cos\vartheta_3&\cos\vartheta_1\\
 \cos\vartheta_2&\cos\vartheta_3&1&
                                      \cos(\vartheta_1+\vartheta_3+jh)\\
 \cos(\vartheta_4-h)&\cos\vartheta_1&
                    \cos(\vartheta_1+\vartheta_3+jh)&1
 \end{pmatrix}.
 \label{eq:screenedangleGram}
\end{equation}
Calling this a formal Gram matrix does not assert that the
angle data are realized by a convex tetrahedron.

At $h=0$ it is the Euclidean planar Gram matrix
\begin{equation}
 (\mathbf G_0)_{ab}=\cos(\phi_a-\phi_b),\qquad
 (\phi_1,\phi_2,\phi_3,\phi_4)=
 (0,\vartheta_2+\vartheta_3,\vartheta_2,
                       \vartheta_1+\vartheta_2+\vartheta_3).
 \label{eq:screenedplanarGram}
\end{equation}
Since $0<\vartheta_2<\pi$, these unit vectors span a plane.
Therefore
\begin{equation}
 \boxed{\qquad \operatorname{rank}\mathbf G_0=2.\qquad}
 \label{eq:screenedranktwo}
\end{equation}
For equal half-weights it is
\[
 \mathbf G_0=
 \begin{pmatrix}1&-1&0&0\\-1&1&0&0\\
                 0&0&1&-1\\0&0&-1&1\end{pmatrix}.
\]
In this equal-weight example the vectors $(1,1,0,0)$ and
$(0,0,1,1)$ lie in the kernel, and the two remaining eigenvalues
are $2$. For comparison, a regular Euclidean tetrahedron has
normal Gram matrix with diagonal $1$ and off-diagonal $-1/3$:
its eigenvalues are $0,4/3,4/3,4/3$. Thus a vanishing determinant
alone does not distinguish a planar configuration from a genuine
Euclidean tetrahedron; the number of independent normal directions does.

By contrast, the outward normals of a nondegenerate Euclidean
tetrahedron span three dimensions. This difference holds for
every four-weight choice on the screening locus, and remains
under varying such weights because rank at most two is closed.

There is already a finite-$b$ obstruction. If $\mathcal C_{ab}$
denotes the cofactor matrix, then
\begin{equation}
 \mathcal C_{1+j,1+j}=0,\qquad
 \mathcal C_{4-k,4-k}=0.
 \label{eq:screenedzerocofactors}
\end{equation}
Each identity is the determinant of three planar unit vectors.
At a finite vertex of a compact hyperbolic tetrahedron, the
three incident face normals lie in the positive-definite
orthogonal complement of a timelike vertex vector and are
independent. Their Gram determinant is positive. Thus
\eqref{eq:screenedzerocofactors} excludes compact tetrahedra
already before taking the limit.

For completeness, the determinant and orientation checks are
also explicit. Each $\mathbf G_{jk}$ is a planar rank-two
Gram matrix perturbed in a single symmetric off-diagonal entry.
The entry, its increment $\epsilon$, and a positive factor $q$
are
\[
\begin{array}{c|c|c|c}
(j,k)&\text{entry}&\epsilon&q\\ \hline
(0,0)&14&\cos(\vartheta_4-h)-\cos\vartheta_4&
                                      \sin^2\vartheta_3\\
(0,1)&13&\cos\vartheta_2-\cos(\vartheta_2+h)&
                                      \sin^2\vartheta_1\\
(1,0)&24&\cos\vartheta_1-\cos(\vartheta_1+h)&
                                      \sin^2\vartheta_2\\
(1,1)&23&\cos\vartheta_3-\cos(\vartheta_3+h)&
                                      \sin^2(\vartheta_4-h)
\end{array}
\]
In every row, $\det\mathbf G=-q\epsilon^2$ and the cofactor
at that off-diagonal entry is $-q\epsilon$. For fixed interior
weights and small positive $h$, $\epsilon>0$. Hence the matrices
have Lorentz signature but fail the positive off-diagonal
cofactor condition for an ordinary convex hyperbolic
tetrahedron; see~\cite[Theorem~2.1]{BellettiYang}.
The zero principal cofactors therefore do not justify calling
these configurations tetrahedra with two ideal vertices.

\subsection{The screened classical saddle also degenerates}

The dilogarithmic potential gives an independent check. At
$b=0$, write $X=A_1,Y=A_2,Z=A_3$. The reflected screening
labels satisfy
\[
 A_4=XYZ,\qquad A_s=-XY,\qquad A_t=-YZ.
\]
Set
$P=X^2Y^2=e^{2\pi i\delta_0}$ and
$R=Y^2Z^2=e^{2\pi i\delta_1}$.
Substitution into~\cite[eq.~(3.1)]{TeschnerVartanov} cancels
two pairs of dilogarithms and leaves
\begin{equation}
 U(u)=\operatorname{Li}_2(u)+\operatorname{Li}_2(PRu)
       -\operatorname{Li}_2(Pu)-\operatorname{Li}_2(Ru).
 \label{eq:screenedclassicalpotential}
\end{equation}
The exponentiated saddle equation $uU_u=-2\pi i$ becomes
\[
 (1-Pu)(1-Ru)=(1-u)(1-PRu),
 \qquad (1-P)(1-R)u=0.
\]
For nonbalanced weights there is no nonzero regular saddle.
Before cancellation, the equation contains the factor
$(1-(XYZ)^2u)^2$; its double root is at the very branch
factors that cancel in~\eqref{eq:screenedclassicalpotential}.
This is not a regular double critical point of the reduced
phase. If $P=1$ or $R=1$, the reduced potential itself vanishes
on a consistent local branch. Thus the finite screening
operator survives in a singular pole/branch collision regime,
rather than in the ordinary compact Euclidean saddle regime.

\subsection{A genuine compact family: the regular tetrahedron}

For comparison, prescribe the actual six interior dihedral
angles $\eta_e$ of a small compact tetrahedron and set
\begin{equation}
 \alpha_e=\frac{\pi-\eta_e}{2\pi b}.
 \label{eq:actualtetrahedralcharges}
\end{equation}
These quantum labels encode edge dihedrals. The $\mu_i$ of the
flat boundary metric encode intrinsic vertex deficits, so the
two kinds of data must be related geometrically before they
can be identified.

The compact family is a separate geometric test. It uses six
edge dihedral angles, whereas the preceding screening calculation
uses four vertex-deficit weights and two intermediate labels.
The distinction can already be seen on a regular Euclidean
tetrahedron: each face angle is $\pi/3$, each intrinsic vertex
opening is $\pi$, and each interior edge dihedral is
$\arccos(1/3)$. These three angles describe different features.

For a regular tetrahedron of edge length $\ell$, let $\beta$
be a face angle and $\eta$ an interior dihedral. Hyperbolic
trigonometry gives
\begin{equation}
 \cos\beta=\frac{\cosh\ell}{1+\cosh\ell},\qquad
 \cos\eta=\frac{\cosh\ell}{1+2\cosh\ell},\qquad
 A=4\pi-12\beta.
 \label{eq:regulartetrahedrongeometry}
\end{equation}
At the Euclidean limit
$\eta_{\mathrm E}=\arccos(1/3)$. Its angle Gram has
\[
 \det G=(1+\cos\eta)^3(1-3\cos\eta),
\]
and rank three at $\eta_{\mathrm E}$.
All six limiting heavy edge weights are
\[
 2b\alpha_e=1-\eta_{\mathrm E}/\pi\simeq0.608173.
\]
In the reflected screening convention,
$2bd_{\mathrm{out}}\to2-\mu_4>1$, whereas this geometric
prescription gives the same number below one for all six
labels. Moreover the screened internal charges are sums of
pairs of external charges. The rank obstruction is independent
of reflected representatives; these explicit prescriptions
already differ at leading order $1/b$, beyond an order-$b$
normalization shift.

Expansion of~\eqref{eq:regulartetrahedrongeometry} gives
\begin{align}
 A&=\sqrt3\,\ell^2-\frac{\sqrt3}{8}\ell^4+O(\ell^6),
 \notag\\
 t:=\eta_{\mathrm E}-\eta
   &=\frac{\ell^2}{12\sqrt2}
       -\frac{23\ell^4}{1152\sqrt2}+O(\ell^6)
     =\frac{A}{12\sqrt6}+O(A^2).
 \label{eq:regularareadihedral}
\end{align}
Using $\dd W=-3\ell\,\dd\eta$ and $W(0)=0$,
\begin{equation}
 \boxed{\quad
 W(A)=\frac{A^{3/2}}{6\sqrt2\,3^{3/4}}
       \left(1-\frac{A}{10\sqrt3}+O(A^2)\right).
 \quad}
 \label{eq:regularvolumebyarea}
\end{equation}
In edge length this is
$W=\ell^3/(6\sqrt2)-23\ell^5/(480\sqrt2)+O(\ell^7)$.

\subsection{The regular classical fold and its scale}

Put all six exponentiated angles equal to $a=e^{i\eta}$.
The TV potential becomes
\[
 U(u,\eta)=\operatorname{Li}_2(u)
          +3\operatorname{Li}_2(a^4u)
          -4\operatorname{Li}_2(-a^3u).
\]
Exponentiating the saddle equation gives
$(1+a^3u)^4=(1-u)(1-a^4u)^3$. After removing
$u(a+1)^2$, the quadratic is
\begin{equation}
 P_\eta(u)=a^8(a^2-2a+3)u^2
              -3a^4(a-1)^2u+(3a^2-2a+1).
 \label{eq:regularsaddlequadratic}
\end{equation}
Its discriminant is
\[
 -a^8(a+1)^2(3a^2-2a+3),\qquad
 3a^2-2a+3=2a(3\cos\eta-1).
\]
Thus the saddles merge simply at the Euclidean angle.
The double root and its real logarithmic lift are
\[
 a_{\mathrm E}=\frac{1+2\sqrt2\,i}{3},\qquad
 u_{\mathrm E}=a_{\mathrm E}^{-5}
   =\frac{241+22\sqrt2\,i}{243},\qquad
 v_{\mathrm E}=2\pi-5\eta_{\mathrm E}.
\]
None of the dilogarithm arguments is $1$ at this point.

A double root of the saddle equation is the point at which the
two separate Gaussian approximations fail: the second derivative
vanishes, so their inverse square-root Hessians would diverge.
To obtain a useful approximation one must keep the next nonzero
term, the cubic, and vary the parameter at the same time. The
derivative calculation below both checks that this is a simple
fold and fixes its constants; a double algebraic root alone would
not check the analytic phase or its branch singularities.

For $u=e^{iv}$ the action, including a term independent of
$v$, is $S(v,\eta)=U(e^{iv},\eta)-2\pi v+\text{const.}$
Choose the local branches satisfying the saddle equation.
Direct differentiation at the fold yields
\[
 S_v=S_{vv}=0,\qquad
 S_{vvv}=54i,\qquad S_{v\eta}=18\sqrt2\,i.
\]
For example,
\[
 S_{v\eta}=-12\left(\frac{a^3u}{1+a^3u}
                         +\frac{a^4u}{1-a^4u}\right).
\]
Writing $x=v-v_{\mathrm E}$ gives the concrete normal form
\begin{equation}
 \boxed{\quad
 S(v_{\mathrm E}+x,\eta_{\mathrm E}-t)
       -S(v_{\mathrm E},\eta_{\mathrm E}-t)
 =i(9x^3-18\sqrt2\,t x)
       +O(x^4+t x^2+t^2|x|).
 \quad}
 \label{eq:regularcubicfold}
\end{equation}

Use the saddle exponential preceding
\cite[eq.~(3.11)]{TeschnerVartanov}, namely $e^{S/(ih)}$
with $h=2\pi b^2$. Rescaling $x=(h/27)^{1/3}y$ gives
$y^3/3-\zeta y$ as the leading exponent, where
\begin{equation}
 \boxed{\quad
 \zeta=\frac{6\sqrt2\,t}{h^{2/3}}
       =\frac{A}{2\sqrt3\,h^{2/3}}
          +O\left(\frac{A^2}{h^{2/3}}\right).
 \quad}
 \label{eq:regularAiryvariable}
\end{equation}
This proves the cubic normal form of the \emph{classical}
phase and derives the scale $A\asymp b^{4/3}$ with its
normalization. A smooth-amplitude integral on a specified
Airy contour has the corresponding uniform asymptotic; see
\cite[Section~36.12]{DLMF}. Section~\ref{sec:quantumAiry}
studies the corresponding assertion for the exact regular-family
integral with an explicit scalar regularization. The local
classical calculation here does not establish the required
uniform contour and tail estimates for that quantum integral.

With TV's choice $\Delta-5\pi^2/3$ for the term independent of
$v$, the volume formula of~\cite{MurakamiYano}, in TV's doubled
potential convention, gives
$\operatorname{Im}S_\pm=\pm2W$. Thus the two saddle
exponentials contain $e^{\pm2W/h}$. If the volume exponent
is written as $W/\hbar_{\mathrm{vol}}$, its parameter is
\begin{equation}
 \hbar_{\mathrm{vol}}=\pi b^2=h/2.
 \label{eq:TVvolumehbar}
\end{equation}
The leading cubic saddle values agree:
$(2/3)\zeta^{3/2}\sim2W/h$.

\subsection{Correction of the integration-coordinate dictionary}

An earlier version of this monograph placed a dense pole lattice at the
classical fold. That conclusion was incorrect: the logarithmic coordinate
of the dilogarithm potential was matched to the wrong shifted integration
variable. The exact integral fixes the dictionary without an asymptotic
choice of branch.

With all six charges $\alpha=(\pi-\eta)/(2\pi b)$, write the
original integration variable in~\cite[eq.~(2.27)]{TeschnerVartanov}
as $w$. Its integrand is
\begin{equation}
 I_b(w,\eta)=S_b(w-3\alpha)^4
             S_b(4\alpha-w)^3S_b(2Q-w).
 \label{eq:regularquantumintegrand}
\end{equation}
Set $a=e^{i\eta}$. Directly,
\[
 e^{2\pi ib(w-3\alpha)}=-a^3e^{2\pi ibw}.
\]
Thus the variable $u$ in the classical potential is
\begin{equation}
 u=e^{2\pi ibw},\qquad v=2\pi bw.
 \label{eq:correctquantumcoordinate}
\end{equation}
The $Q/2$-shift used previously is incompatible with this
identification. A second check is the exact shift quotient
\begin{align}
 \frac{I_b(w+b,\eta)}{I_b(w,\eta)}
 &=
 \frac{[2\sin\pi b(w-3\alpha)]^4}
 {[2\sin\pi b(4\alpha-w-b)]^3
  [2\sin\pi b(2Q-w-b)]}\notag\\
 &\longrightarrow
 \frac{(1+a^3u)^4}{(1-a^4u)^3(1-u)}.
 \label{eq:exactintegrandquotient}
\end{align}
The limit reproduces the saddle polynomial
\eqref{eq:regularsaddlequadratic}.

The lift of the double root lying in the separating contour strip is
\[
 v_*=4\pi-5\eta_{\mathrm E},\qquad
 w_*=\frac{v_*}{2\pi b}.
\]
It differs by $2\pi$ from the logarithmic lift used in the classical
calculation; all its local derivatives remain the same. At this lift,
\begin{equation}
 \begin{aligned}
 w_*-3\alpha_{\mathrm E}
       &=\frac{\pi-2\eta_{\mathrm E}}{2\pi b},\\
 4\alpha_{\mathrm E}-w_*&=\frac{\eta_{\mathrm E}}{2\pi b},\\
 2Q-w_*&=\frac{5\eta_{\mathrm E}}{2\pi b}+2b.
 \end{aligned}
 \label{eq:correctfoldarguments}
\end{equation}
All three arguments lie in $0<\operatorname{Re}z<Q$ for
sufficiently small $b$, a fixed scaled distance from its boundary.
There is therefore no integration-pole accumulation at the fold.
The external trivalent factors do have parameter poles; their
regularization and the further analytic requirements for a uniform
transition are treated in the next section. Pole separation in
this local strip alone does not control the complete integration
contour or its tails.

\paragraph{Computational companion.}
The reproducible checks are described in
\path{experiments/quantum_surgery/README.md}, with results in
\path{artifacts/quantum_surgery/result.json}. Exact finite algebra
checks support the stated character and Gram identities;
multiprecision comparisons independently evaluate the finite
period formulas and Euler arcs. These checks supply neither
the infinite-spectral estimate~\eqref{eq:screeningtailinput}
nor a global quantum Airy certificate.

